\chapter{Introduction}
\label{chap:introduction}

\ac{ICS} are considered to be critical systems as they are responsible for monitoring and controlling important infrastructures such as water systems, power plants, and nuclear plants. Typically, \ac{ICS} consist of two major areas which are, \ac{IT}, and \ac{OT}. \ac{IT} focuses more on the standard computer interconnections via traditional communication protocols such as \ac{HTTP}, and \ac{HTTPs}. while \ac{OT} focus more on the operational level which includes special hardware and software for monitoring and controlling sensors and actuators. \ac{OT} hierarchy usually consists of field, control, and supervisory layers. Each layer is responsible for certain tasks which works as pipeline. Field layer includes sensors and actuators, while control layer consists of number of \ac{PLCs} which are responsible for managing and controlling all field layer devices. Sensors measure the current state of the underlying physical process, then they report it to \ac{PLCs}. \ac{PLCs} implement certain control logic which is used to evaluate the incoming sensor readings and take decisions based on them. Based on these decisions, \ac{PLCs} may send certain commands to the field actuators to adjust the current state of the physical process. In the supervisory layer, there exist Supervisory Control and Data Acquisiontion (SCADA) system performs a high-level visualization and control over the control and field layers and provides an overview of the current process state to the working engineer (operator). \ac{OT} used to be \textit{air-gapped}, which means that they were designed to be isolated from external enviroments, and they were secure via \textit{obsecurity}. Currently industrial organizations are increasingly connecting their \ac{OT} network with the corporate networks to improve business and operational efficiency, but this comes at a cost. \ac{ICS} are exposed to cyber-physical attacks such as security breaches in cyberspace that adversely affect the physical process. To defend \ac{ICS} from such attacks, \ac{IT} is important to adopt multi-layered approach. First, \ac{IT} is desirable to block cyber attacks before they reach the \ac{OT} network. Indeed, most attacks against \ac{ICS} start by gaining access to the business network and then laterally moving to the \ac{OT} network. Second, \ac{IT} is important to have anomaly detection tools that works as outlier detection, which are responsible for identification and observation of events that deviate from the normal behavior, and report such cases to the working operators to handle them. Anomaly detection has a long history in the field of statistics, where analysts and scientists would study charts looking for any element that appeared abnormal. Today, anomaly detection leverages artificial intelligence and machine learning to automatically identify unexpected changes in a data set's normal behavior. \ac{GANs} are a machine learning framework where two neural networks - a generator and a discriminator - compete against each other in a zero-sum game to produce realistic, synthetic data. The generator creates fake data, while the discriminator evaluates \ac{IT} against real data, and the main goal is to teach the discriminator how \ac{IT} can differentiate normal data from malicious data. Major problem in \ac{ICS} is that some of the used protocols are vendor specific and not open-source, which prevents many researchers and developers from building anomaly detection tools which are able to understand the protocol and allow defenders to detect attacks. One solution to bypass such problem is to develop reverse engineering methods in a trial to understand the structure of the protocol, and determine the role of each chunk. The goal of this study project is to implement and evaluate different approaches that are able to extract expressive patterns, i.e., features, from network traffic collected in the \ac{OT} area of a given \ac{ICS} via reverse engineering approach and build different \ac{ML} and \ac{DL}  models to work as classifiers. We assume that the \ac{ICS} operator utilizes proprietary binary protocols, and the approaches to be developed can not directly parse the content of the transmitted network packets. We also rely on performing reverse engineering method to extract the physical reading out from the network packets. Then we rely on different \ac{ML} and \ac{DL} methods to train a classifier with the extracted chunks in order to create a classification model. Finally this model is to be used as a anomaly detection tool to identify possible abnormal operation in \ac{ICS}.  In our analysis we focus on \ac{S7comm} protocol which is one of the leading protocols used in industrial networks. Siemens S7 \ac{PLCs} are estimated to consititue over 30\% of the world-wide PLC market. The main reason that we choose \ac{S7comm} is its proprietary nature and it was not publicly available before it was reverse engineered. So we used 2017QUTS7Comm dataset to perform all our analysis. 

{\color{red} todo: Research Objective}

{\color{gray} Main Question:
Can anomaly detection methods be designed to operate without relying on protocol knowledge and achieve good performance on reverse engineered traffic?}
